\sectionkicker{Source-backed technical notes}
\subsection*{フロンティア競争は性能表から長文脈・マルチモーダルなインタラクションへ: Technical Notes}
本欄は記事本文で圧縮した一次資料上の情報を、比較・再検証しやすい形で整理したものである。時系列、確認済みの主張、留意点、一次資料URLを掲載する。
\smallskip
\noindent{\small\textit{Technical Notesでは、一次資料から確認した公開・提供・研究上の事実を記載する。数値・能力評価や提供元・著者による比較表現は、独立再現済みの結果ではなく帰属claimとして扱う。}}\par
\medskip
\subsection*{Theme at a glance}
\addcontentsline{toc}{subsection}{Theme at a glance}
\begin{center}
\footnotesize
\begin{tabularx}{\linewidth}{@{}p{0.20\linewidth}p{0.10\linewidth}p{0.14\linewidth}X@{}}
\toprule
資料 & 位置づけ & 種別 & 時系列 \\
\midrule
GPT-4 / GPT-4V 2023 lifecycle & 主要資料 & モデル更新 & — \\
Claude 2023 model family & 主要資料 & モデル更新 & — \\
Gemini 2023 launch family & 主要資料 & モデル & — \\
PaLM 2 & 補足資料 & モデル & — \\
Lost in the Middle & 補足資料 & 論文 & — \\
\bottomrule
\end{tabularx}
\end{center}
\normalsize

\begin{technicalnote}{GPT-4 / GPT-4V 2023 lifecycle}{主要資料}
\begingroup
% reader-facing Technical Notes break policy
\widowpenalty=10000
\clubpenalty=10000
\displaywidowpenalty=10000
\begin{tabularx}{\linewidth}{@{}>{\bfseries}p{0.22\linewidth}X@{}}
組織 & OpenAI \\
種別 & モデル更新 \\
時系列 & — \\
\end{tabularx}
\smallskip
{\bfseries 一次資料から整理したtechnical points}
\begin{itemize}[leftmargin=1.5em,itemsep=0.35em]
\item \textbf{一次情報で確認できる事実}: OpenAIの選定済み一次資料では、対象event近傍の一次資料から factuality evaluation を確認できる。
\end{itemize}
\begin{minipage}{\linewidth}
% reader-facing Technical Notes generic boundary/source tail group
\begin{samepage}
{\bfseries 一次資料}
\begingroup\sloppy
\Urlmuskip=0mu plus 2mu\relax
\begin{itemize}[leftmargin=1.5em,itemsep=0.25em]
\item {\scriptsize\url{http://arxiv.org/abs/2303.08774v6}}
\item {\scriptsize\url{https://openai.com/index/chatgpt-can-now-see-hear-and-speak}}
\item {\scriptsize\url{https://openai.com/index/gpt-4-research}}
\item {\scriptsize\url{https://openai.com/index/gpt-4v-system-card}}
\end{itemize}
\endgroup
\end{samepage}
\end{minipage}
\endgroup
\end{technicalnote}
\medskip

\begin{technicalnote}{Claude 2023 model family}{主要資料}
\begingroup
% reader-facing Technical Notes break policy
\widowpenalty=10000
\clubpenalty=10000
\displaywidowpenalty=10000
\begin{tabularx}{\linewidth}{@{}>{\bfseries}p{0.22\linewidth}X@{}}
組織 & Anthropic \\
種別 & モデル更新 \\
時系列 & — \\
\end{tabularx}
\smallskip
{\bfseries 一次資料から整理したtechnical points}
\begin{itemize}[leftmargin=1.5em,itemsep=0.35em]
\item \textbf{一次情報で確認できる事実}: Anthropicの選定済み一次資料では、2023年3月の初回公開では Claude と軽量・低価格・高速版の Claude Instant の2版が提示され、Claude は chat interface と developer console の API から利用可能とされた。 2023年5月には Claude の context window が 9K から 100K tokens へ拡張され、100K context は API で利用可能とされた。 2023年7月の Claude 2 は API と claude.ai の public-facing beta で提供され、1 prompt あたり最大 100K tokens の入力をサポートした。 2023年11月の Claude 2.1 は API と claude.ai で提供され、200K-token context window を導入し、developer-defined functions / APIs と連携する tool use を beta feature として追加した。
\end{itemize}
\begin{minipage}{\linewidth}
% reader-facing Technical Notes generic boundary/source tail group
\begin{samepage}
{\bfseries 一次資料}
\begingroup\sloppy
\Urlmuskip=0mu plus 2mu\relax
\begin{itemize}[leftmargin=1.5em,itemsep=0.25em]
\item {\scriptsize\url{https://www.anthropic.com/news/100k-context-windows}}
\item {\scriptsize\url{https://www.anthropic.com/news/claude-2}}
\item {\scriptsize\url{https://www.anthropic.com/news/claude-2-1}}
\item {\scriptsize\url{https://www.anthropic.com/news/introducing-claude}}
\end{itemize}
\endgroup
\end{samepage}
\end{minipage}
\endgroup
\end{technicalnote}
\medskip

\begin{technicalnote}{Gemini 2023 launch family}{主要資料}
\begingroup
% reader-facing Technical Notes break policy
\widowpenalty=10000
\clubpenalty=10000
\displaywidowpenalty=10000
\begin{tabularx}{\linewidth}{@{}>{\bfseries}p{0.22\linewidth}X@{}}
組織 & Google \\
種別 & モデル \\
時系列 & — \\
\end{tabularx}
\smallskip
{\bfseries 一次資料から整理したtechnical points}
\begin{itemize}[leftmargin=1.5em,itemsep=0.35em]
\item \textbf{一次情報で確認できる事実}: Googleの選定済み一次資料では、Gemini 1.0 は Ultra・Pro・Nano の3サイズとして発表され、モデル自体は text・code・audio・image・video を横断する native multimodality を前提に設計された。 2023年12月13日時点で developer / enterprise 向けに公開されたのは Gemini Pro 系で、Gemini API の text版は 32K context で text input / text output、Gemini Pro Vision endpoint は text と imagery を入力して text を出力する構成だった。 Gemini Ultra は2023年12月の API 公開時点では追加の fine-tuning と safety testing 後に翌年提供予定とされており、2023年中の Gemini Pro availability と同一視しない。 Gemini Nano は Pixel 8 Pro 上の on-device model として2023年12月6日に展開され、Recorder の Summarize を端末上で実行し、Gboard Smart Reply は developer preview として段階展開された。
\end{itemize}
\begin{minipage}{\linewidth}
% reader-facing Technical Notes generic boundary/source tail group
\begin{samepage}
{\bfseries 一次資料}
\begingroup\sloppy
\Urlmuskip=0mu plus 2mu\relax
\begin{itemize}[leftmargin=1.5em,itemsep=0.25em]
\item {\scriptsize\url{http://arxiv.org/abs/2312.11805v5}}
\item {\scriptsize\url{https://blog.google/innovation-and-ai/products/gemini-api-developers-cloud/}}
\item {\scriptsize\url{https://blog.google/innovation-and-ai/technology/ai/google-gemini-ai/}}
\item {\scriptsize\url{https://blog.google/products-and-platforms/devices/pixel/pixel-feature-drop-december-2023/}}
\end{itemize}
\endgroup
\end{samepage}
\end{minipage}
\endgroup
\end{technicalnote}
\medskip

\begin{technicalnote}{PaLM 2}{補足資料}
\begingroup
% reader-facing Technical Notes break policy
\widowpenalty=10000
\clubpenalty=10000
\displaywidowpenalty=10000
\begin{tabularx}{\linewidth}{@{}>{\bfseries}p{0.22\linewidth}X@{}}
組織 & Google \\
種別 & モデル \\
時系列 & — \\
\end{tabularx}
\smallskip
{\bfseries 一次資料から整理したtechnical points}
\begin{itemize}[leftmargin=1.5em,itemsep=0.35em]
\item \textbf{一次情報で確認できる事実}: Googleの選定済み一次資料では、PaLM 2 は2023年5月10日に次世代 language model として公開され、Google は multilingual・reasoning・coding を主要な能力軸として説明し、100超の言語を含む multilingual text や公開 source-code dataset を学習面の特徴として挙げた。 PaLM 2 family は小さい順に Gecko・Otter・Bison・Unicorn の4サイズとして提示され、Gecko は mobile device 上で offline でも対話用途に使える軽量 variant として位置づけられた。 公開時点で PaLM 2 は25超の Google products / features に組み込まれ、developer は PaLM 2 model への sign-up、enterprise customer は Vertex AI からの利用という別の availability boundary が示された。 Med-PaLM 2 や Sec-PaLM は PaLM 2 family の用途特化 variant として記述されており、Med-PaLM 2 の限定 Cloud customer feedback や Sec-PaLM の Google Cloud availability を、base PaLM 2 全体の提供条件として一般化しない。
\end{itemize}
\begin{minipage}{\linewidth}
% reader-facing Technical Notes generic boundary/source tail group
\begin{samepage}
{\bfseries 一次資料}
\begingroup\sloppy
\Urlmuskip=0mu plus 2mu\relax
\begin{itemize}[leftmargin=1.5em,itemsep=0.25em]
\item {\scriptsize\url{https://blog.google/innovation-and-ai/products/google-palm-2-ai-large-language-model/}}
\end{itemize}
\endgroup
\end{samepage}
\end{minipage}
\endgroup
\end{technicalnote}
\medskip

\begin{technicalnote}{Lost in the Middle}{補足資料}
\begingroup
% reader-facing Technical Notes break policy
\widowpenalty=10000
\clubpenalty=10000
\displaywidowpenalty=10000
\begin{tabularx}{\linewidth}{@{}>{\bfseries}p{0.22\linewidth}X@{}}
組織 & authors \\
種別 & 論文 \\
時系列 & — \\
\end{tabularx}
\smallskip
{\bfseries 一次資料から整理したtechnical points}
\begin{itemize}[leftmargin=1.5em,itemsep=0.35em]
\item \textbf{一次情報で確認できる事実}: authorsの選定済み一次資料では、著者らは multi-document question answering と key-value retrieval の2タスクを用い、long context 内で relevant information の位置を変えたときの利用性能を評価した。 著者実験では relevant information が入力の先頭または末尾にある場合に性能が高くなりやすく、中間に置かれた場合に大きく低下する傾向が、明示的な long-context model を含めて観測された。 したがって公称 context-window length は、その全域の情報を一様に有効利用できることの証拠にはならない。この点は論文著者の実験結果として扱い、個別モデルの全条件へ一般化しない。
\end{itemize}
\begin{minipage}{\linewidth}
% reader-facing Technical Notes generic boundary/source tail group
\begin{samepage}
{\bfseries 一次資料}
\begingroup\sloppy
\Urlmuskip=0mu plus 2mu\relax
\begin{itemize}[leftmargin=1.5em,itemsep=0.25em]
\item {\scriptsize\url{http://arxiv.org/abs/2307.03172v3}}
\end{itemize}
\endgroup
\end{samepage}
\end{minipage}
\endgroup
\end{technicalnote}
\medskip
